%% 自動生成（gen.py）
\bktocchap{5-1}{平均の速さと瞬間の速さ}{2時間で100 km走った車は、ずっと時速50 km？}{1}{2}{2}{18}{C}
\bktocchap{5-2}{時間をどんどん短くしてみよう}{0.1秒、0.01秒…と縮めると、ある値に近づく}{1}{2}{2}{40}{C}
\bktocchap{5-3}{かぎりなく近づく：極限と接線}{曲線にぴったりふれる直線の傾き}{1}{3}{3}{60}{C}
\bktocchap{5-4}{表から導関数を見つけよう}{表をながめると、微分のきまりが見えてくる}{1}{3}{3}{78}{C}
\bktocchap{5-5}{位置→速度→加速度}{2回微分すると、何がわかる？}{1}{3}{3}{104}{C}
\bktocchap{5-6}{微分の道具箱：積と合成関数}{歯車をつなぐと、回る速さはかけ算になる}{1}{4}{4}{126}{C}
\bktocchap{5-7}{いちばん高いのはどこ？　最大と最小}{いちばん高い点では、速度が0になる}{1}{3}{3}{154}{C}
\bktocchap{5-8}{短い時間なら、かけ算で近似できる}{「次の位置」を少しずつ予想する}{1}{2}{2}{174}{C}
